\section{一个 Pipe 怎样让 reader 等待数据}

创建一条空 Pipe，让 reader 先执行 \code{read}。此时没有字节可以复制，
writer 又仍然存活，所以返回 EOF 是错误的；持续轮询也会让 reader 占满 CPU。
内核先用 \code{TryRead} 得到 \code{WouldBlock}，随后让当前 Thread 睡眠。
writer 写入数据后把 reader 移回 Ready，reader 再次运行时从头检查条件。

\begin{pdfdiagram}
  \node[os hardware] (read) {read(fd)};
  \node[os block,right=of read] (try) {TryRead\\缓冲为空};
  \node[os state,right=of try] (wait) {登记 reader\\Thread=Blocked};
  \node[os block,below=of wait] (write) {writer 写入\\唤醒一个 reader};
  \node[os state,left=of write] (ready) {Thread=Ready};
  \node[os hardware,left=of ready] (retry) {重新 TryRead\\复制字节};
  \draw[os arrow] (read) -- (try);
  \draw[os arrow] (try) -- node[above]{WouldBlock} (wait);
  \draw[os arrow] (wait) -- (write);
  \draw[os arrow] (write) -- (ready);
  \draw[os arrow] (ready) -- (retry);
\end{pdfdiagram}
\diagramnote{唤醒只改变调度资格，不把某一段数据预留给 reader；恢复运行后仍要重新检查缓冲区。}

\topic{先把 read 和 write 的四种结果列清楚}
Pipe 是没有消息边界的字节流。一次 write 可以被多次短 read 取走，一次 read
也可能合并多次 write 的内容。缓冲区是否为空、是否已满，以及两类端点是否
仍然存活，共同决定本次操作：

\begin{center}
\begin{osgridtabular}{L{0.32\linewidth}L{0.27\linewidth}L{0.28\linewidth}}
当前状态 & read & write \\ \hline
空且 writer 存活 & WouldBlock & 可以写入 \\ \hline
存在数据 & 返回当前可用的短读 & 取决于剩余空间 \\ \hline
空且 writer 全关闭 & EOF & 不适用 \\ \hline
reader 全关闭 & 不适用 & BrokenPipe \\ \hline
满且 reader 存活 & 可以读取 & WouldBlock \\ \hline
\end{osgridtabular}
\end{center}

EOF 不是缓冲区里的特殊字节。writer 关闭时若仍有数据，reader 必须先把数据
读完；缓冲区变空后的下一次 read 才返回 EOF。BrokenPipe 也要等最后一个
reader 引用消失才能成立，关闭其中一个复制出来的 fd 还不够。

\topic{锁内复查堵住丢失唤醒的窗口}
reader 第一次检查为空后，writer 可能立刻写入。若 reader 接下来无条件入队，
writer 的唤醒已经过去，它会永久留在 Blocked。实现把以下动作放进同一临界区：

\begin{enumerate}[label=\arabic*.]
  \item 取得 Pipe 锁，再检查 buffered 与 writer 数量；
  \item 条件已经改变时立即返回，不进入等待队列；
  \item 条件仍未成立时，把当前 Thread 链到 reader 队尾；
  \item 在调度器中把该 Thread 从 Running 改为 Blocked；
  \item 释放锁并切换到另一个 Ready Thread。
\end{enumerate}

writer 在同一把锁下增加 buffered，然后从 reader 队首取出一个 Thread，
将它改为 Ready。超时或信号也可能同时尝试完成这次等待，因此每个等待记录只
允许一个完成原因成功写入。先把具体竞争关进锁和单赢家状态后，我们才把这条
规则称为等待协议的不变量。

\topic{64 KiB 容量按需要取得 4 KiB 页面}
普通用户 Pipe 的逻辑容量是 \(65536\) 字节，分成
\[
  \frac{65536}{4096}=16
\]
个后备页槽。刚创建时，buffered 和 allocated page count 都为零。写入某个
页覆盖的区间时才经 allocator callback 取得物理 frame，并通过 direct-map
访问它。这样，1024 条空 Pipe 不会仅因声明容量就预占 64 MiB。

环形状态保存 read index、write index 和 buffered byte count。索引相等时，
buffered 区分空与满；可写空间始终为
\[
  \mathrm{writable}=65536-\mathrm{buffered}.
\]
逻辑位置 \(i\) 落到
\[
  \mathrm{page}=\left\lfloor\frac{i}{4096}\right\rfloor,\qquad
  \mathrm{offset}=i\bmod4096.
\]

Pipe 不读取 E820，也不操作 buddy 或页表层级。它只调用固定签名的页面
acquire/release 回调；宿主测试给它一组普通内存，目标 Kernel 则把回调接到
真实物理页分配器。

\topic{第三个页面申请失败时，前两个也要归还}
一次跨页 write 可能需要三个尚未分配的页面。若前两次申请成功、第三次失败，
代码不能推进 write index，也不能把前两页留在 Pipe 中。实现先计算缺失页，
把本次新取得的地址暂存在局部数组中；所有页面都准备好以后才复制字节并更新
buffered。中途失败就逆序归还局部数组中的页面。

释放时采用相反的谨慎顺序：release callback 成功以后才清除页槽。若先把地址
清零，回调失败后系统就失去了仍被自己占用的 frame。一次操作要么让所有状态
一起前进，要么恢复到调用前的状态；这种由具体回滚动作保证的性质称为失败
原子性。

\topic{fd 为什么不能直接保存 Pipe 槽号}
现在让同一条 Pipe 的读端出现在两个 fd 中。关闭其中一个 fd 后，另一个仍然
应该能读；fork 以后，父子进程也可能继续引用同一读端。若 fd 直接等于 Pipe
数组下标，内核无法区分“关闭一个名字”和“关闭最后一个读端”。

\begin{center}
\begin{osgridtabular}{L{0.22\linewidth}L{0.25\linewidth}L{0.40\linewidth}}
层次 & 保存位置 & 具体内容 \\ \hline
fd number & 单个 Process & FileTable 中的局部名字与 fd flags \\ \hline
FileTable entry & 单个 FileTable & 指向 FileDescription 的强引用 \\ \hline
FileDescription & KernelObject manager & 读/写能力、共享状态与 finalizer \\ \hline
Pipe & PipeManager slot & 字节、索引、端点数量、等待队列与后备页 \\ \hline
\end{osgridtabular}
\end{center}

\code{dup2(3,0)} 增加的是对同一 FileDescription 的引用，不复制 Pipe
数据。只有最后一个强引用消失，finalizer 才通知 PipeManager 关闭逻辑端点。
这些对象各自负责不同的销毁时刻，这就是后文所说的所有权边界。

\topic{PipeManager 怎样创建和销毁一对端点}
PipeManager 最多保存 1024 个槽，按 RAM 档开放 8、128 或 1024 项。创建时
先取得空槽，再依次建立 reader FileDescription、writer FileDescription、
reader fd 和 writer fd。任一步失败，代码都按相反顺序关闭已经取得的对象；
用户最终只看到两个可用 fd，或者一个明确错误。

两端最后关闭后，manager 释放后备页并归还槽。启动自检会先占满当前档位，
确认额外创建被拒绝，再关闭全部对象并重新创建一次。日志只在结束时打印
capacity、active、peak、creation、release 和 rejection 汇总；逐字节打印会
改变本章正在观察的调度时序。

\topic{从一条 Pipe 过渡到两个程序}
Unix 约定 fd 0、1、2 分别表示标准输入、标准输出和标准错误。普通文件、终端
和 Pipe 可以被放在这些相同的小整数后面，程序因而只调用 read/write，不需要
知道另一端究竟连接了谁。

Shell 要运行 \code{cat | wc} 时，先创建 Pipe，再 fork 两个 child。第一个
child 用 \code{dup2(pipe_writer,1)} 把标准输出接到写端；第二个 child 用
\code{dup2(pipe_reader,0)} 把标准输入接到读端。接线完成后，两边都关闭不再
使用的临时 fd，再分别 exec \code{cat} 和 \code{wc}。

FileTable 的 \code{DuplicateTo} 先验证 source 和目标范围，需要时在锁外
准备新 chunk；重新加锁后复查 source，取得新强引用并替换目标 entry。旧目标
引用在解锁后释放，因为 finalizer 可能关闭 Pipe 并唤醒 Thread，不能在
FileTable 锁内执行。source 等于 destination 时直接成功，否则先关闭目标会
把唯一 source 一同删除。

\topic{多级流水线为什么要先全部启动}
对 \(N\) 个 stage，parent 先创建 \(N-1\) 条 Pipe，然后 fork 全部 child。
child \(i\) 把前一条 Pipe 的 reader 接到 fd 0，把后一条 Pipe 的 writer
接到 fd 1；首尾两端只连接自己需要的一侧。parent 完成所有 fork 后关闭临时
端点，最后才逐个 wait。

若 Shell 启动 stage 0 后立刻等待，它写满 64 KiB 就会进入 Blocked，而读取
数据的 stage 1 尚未创建，系统确定性死锁。若 parent 或无关 child 忘记关闭
writer，末端 reader 即使读完全部数据也看不到 EOF。流水线因此先建立所有
消费者，再等待任何一个生产者结束。

输入重定向把只读文件 duplicate 到 fd 0，输出重定向把 create/truncate 文件
duplicate 到 fd 1。修改只发生在 child 的 FileTable 中，所以 parent 的下一
条 prompt 仍写向终端。

\topic{解析失败发生在创建进程以前}
Shell 先把整行转换为 \code{ShellExecutionPlan}。计划包含 513 字节存储、
128 个参数槽和 16 个 stage；argument 保存存储区内的 offset/length，stage
保存参数范围和可选重定向路径。

解析器识别空格、单双引号、反斜杠、\code{|}、\code{<} 和 \code{>}。空
stage、重复重定向、缺失路径、悬空转义、未闭合引号或容量超限都会使计划清零。
在计划成功以前不调用 pipe、fork 或 open，因此一个语法错误不会遗留半组 child
和 fd。

每个 stage 暂限 8 个参数，计划通过 \code{static\_assert} 保持在 4096 字节
内。用户栈 VMA 虽预留 8 MiB，却只允许相邻页面逐步增长；过大的局部对象可能
让函数序言跨过未提交页，所以更大的计划应迁移到 UserHeap。

\topic{中途 fork 失败时先关闭端点，再等待 child}
第 \(k\) 个 fork 失败时，前 \(k\) 个 child 已经可能运行。parent 先关闭自己
保存的全部 Pipe fd，使正在等待的 child 得到 EOF 或 BrokenPipe，然后 wait
已经发布的 child。若直接返回 prompt，系统会遗留 Zombie、FileDescription
引用和 Pipe 页面。

child 接线失败以 126 退出，exec 或程序查找失败以 127 退出。child 不能返回
Shell 主循环，否则会出现第二个交互 Shell 与 parent 争用键盘。资源按照
“临时 fd、FileDescription、Pipe 页面与槽”的顺序逐层释放；每一层只销毁
自己实际取得的对象。

\topic{运行哪些实验可以看见这些状态}
宿主单元测试分别制造空、满、跨页、回绕、EOF、BrokenPipe 和页面申请失败。
动态模型用固定种子执行 100000 次读、写、关闭和唤醒，每一步都检查 buffered、
端点数量、页面数量与队列成员。Shell 解析测试再输入 4096 条任意字节序列，
成功计划必须满足 stage/fd 边界，失败计划必须保持全零。

QEMU 中先运行一个 producer/consumer：consumer 在空 Pipe 上进入 Blocked，
producer 写入后将它唤醒；writer 关闭以后，consumer 排空数据并读到 EOF。
随后输入一条 16 级流水线，同时使用 15 条 Pipe。运行结束时，active Pipe
应为 0，所有 child 均已 wait，动态页面、FileDescription 和 Zombie 数量回到
运行前。测试脚本设置截止时间，所以丢失唤醒或遗留 writer 会以稳定超时暴露，
而不是让终端无限等待。

完整的 Shell 命令、十九个 freestanding 工具、进程组、Ctrl+C 与前后台切换
在本章后半部分继续。本章到这里已经建立一条更小的因果链：条件不满足时 Thread
进入 Blocked，写入或关闭改变条件，等待者回到 Ready 并重新检查。
